% Paper v2 draft: the new Methods subsections (spec v2 §6). Not \input yet; it joins
% sections/methodology.tex in the writing phase, when the corpus and retrieval
% subsections are rewritten for the v2 release. \pending{...} marks values that come
% from runs not yet done; "% cite:" marks references still to add to references.bib.

\subsection{Decision Roles on a Decision Model}\label{sec:decision}

Two steps of the pipeline make a decision without writing any text. The guard decides whether a turn is an attack, off-topic or a medicine question, and the judge decides whether the evidence suffices to answer. When a language model serves these steps, each decision pays for generated tokens and a full generation round trip, and the confidence the model states in words is a weak signal for when to stop searching. % cite: TASR
We therefore serve both steps with a decision model, Jev from TypeSafe, which generates no text and returns a calibrated probability for each yes-or-no question it receives. % cite: Jev technical report
The steps that write text, namely rephrase, refine and answer, stay on the language model.

Each call to the decision model sends a state and a battery of questions. The state is a JSON object that holds the Vietnamese content of the turn, and every question is an English instruction with a short description of what counts as yes and what counts as no. The batteries are frozen in code, so every run asks the same questions in the same words. The decision model is a separate provider per pipeline role, with its own endpoint, while the graph, the budgets and the answer format stay unchanged.

The guard sends one request per turn after the regular-expression screen. Its state holds the conversation so far and the latest user turn. The battery asks whether the turn is an attack, whether it asks about a medicine and whether it is only small talk or a question about the assistant. It also asks, for each of the 17 intents of the benchmark, whether the turn asks for that kind of information, for example a dose or a contraindication. A turn is blocked when the attack probability reaches 0.70, and it is redirected when both the medicine probability and the small-talk probability stay below 0.5. The intents whose probability reaches 0.5 form the list of asked intents, which the judge and the refine step reuse. The block threshold starts from the value in the vendor's guardrail recipe and is tuned on held-out data, as described below.

The judge sends one request per search round. Its state holds the original question, the standalone query, the asked intents and the evidence as a list of numbered items, each with its source heading and chunk text. For every asked intent, the battery asks whether the evidence covers it for the drug and the population in the question. Three further questions ask whether the evidence suffices overall, whether two items conflict and whether the evidence concerns a different drug. The judge answers when the probability of sufficient evidence and the coverage probability of every asked intent all reach a threshold $\tau$. Otherwise it asks for another search round, and the intents below $\tau$ become the gaps. The probability of sufficient evidence is kept as the judge's confidence. When the guard verdict did not come from the decision model, the judge first asks the intent questions itself, once per turn.

With intents as gaps, the refine step no longer needs a language model. The rephrase step already extracts the drugs or active ingredients that the turn asks about. The template refine step pairs each of these subjects with a fixed Vietnamese phrase for each uncovered intent, for example \textit{liều dùng và cách dùng} for a dose, and keeps at most three new queries, the same limit as the language-model refine step. A judge--refine loop on the decision model with template refine therefore costs no language-model call at all.

The thresholds are never tuned on the end-to-end golden set. They are chosen on a held-out calibration set of 300 benchmark questions, with the evidence that the pipeline retrieves and a label of sufficient evidence derived from the gold sections, together with 100 additional attack and out-of-scope turns for the guard. On this set we report the area under the ROC curve, the expected calibration error and the Brier score of each decision, the decision flip rate over three repeated calls and the effect of reordering the questions. Every raw response is stored with the model version that the service reports, and a run stops if that version changes, because the service does not allow a version to be pinned. The decision model is never the only safety control, since the regular-expression screen runs before it and the answer step keeps its own rules.

This design changes how the cost of a decision grows. A decision is priced by its input tokens only, at \$0.042 per million, and it returns a probability instead of a category. The threshold therefore becomes a parameter that trades evidence quality against search rounds, which Sect.~\ref{sec:gating} exploits.

\subsection{Confidence-Gated Refinement}\label{sec:gating}

A fixed retrieve-more loop spends a judge call and a search round on every turn whose first evidence is already sufficient. We gate the loop instead. The pipeline searches again only when the judge's confidence falls below the threshold $\tau$, so that $\tau=0$ never refines and $\tau=1$ refines until the round limit. We sweep $\tau\in\{0,0.3,0.5,0.7,0.9,1\}$ and report each point as truthfulness against cost and against latency, together with the share of turns that actually refine and the cost of the judge step itself. % cite: Qian et al. 2026 (reporting protocol)
The threshold chosen on Qwen3.5-9B is then applied unchanged to Gemini 3.8 Flash and Claude Opus 5.5, which tests whether it transfers across answer models. As a control, the same sweep runs with the confidence that the language-model judge states in its structured output, which separates the effect of gating from the effect of calibrated probabilities. The selected threshold is \pending{$\tau$ from the sweep on the calibration set}.

\subsection{ReAct Baseline}\label{sec:react}

The pipeline fixes its control flow in code, which raises the question of whether a cheap model would do better if it decided by itself when to search and when to answer. % cite: Anthropic 2024; Ferrazzi et al. 2026; Kumar et al. 2026
To answer this question, we build a ReAct agent that differs from the pipeline only in who controls the loop. It reuses the guard unchanged and then hands the turn to one language model with the conversation history in its messages. There is no rephrase, judge or refine step. The model receives one read-only tool, \texttt{search\_corpus}, which takes a Vietnamese query, runs the retrieval service of the pipeline with the same hybrid retriever and reranker, and returns at most five chunks. Each chunk appears with its heading path and page, as the exact text that the answer step of the pipeline reads, and is numbered in the order in which it was first shown during the turn. The answer cites these numbers under the same rules as the pipeline, and the same citation filter and answer labels apply. The system prompt is short and in Vietnamese. It asks the model to search before answering instead of guessing, notes that most questions need one or two searches, and suggests starting broad and then narrowing. No neighbour-reading tool is offered, because the pipeline has no matching step and each chunk already merges adjacent sections.

The agent shares every budget of the pipeline, which is a condition for a fair comparison. A turn ends when a model response contains no tool call. After six responses with tool calls, or earlier when only one model call remains in the budget or the token budget runs out, the agent adds a note that the search turns are used up and calls the model one last time with tool calls disabled. The evidence budget counts the chunk text placed in the messages. A chunk shown earlier in the turn returns only as its number, and a result that would exceed the budget ends with a note that the evidence budget is spent. A repeated query, compared after case and whitespace normalisation, returns the earlier result with a note instead of searching again. Malformed arguments and failed searches come back as tool messages that start with \texttt{Error}, so that the model can correct itself instead of ending the turn.

The two systems call the model through different interfaces, structured output for the pipeline and native tool calling for the agent, and we report this difference because the interface itself can affect results. % cite: Patel et al. 2026
Everything else is held fixed, namely the corpus, retriever, reranker, chunks, answer format, reasoning setting, temperature and evaluation judge. The agent cannot use the decision model, since every step it takes is a generation. Each run is executed with Qwen3.5-9B, Gemini 3.8 Flash and Claude Opus 5.5 on the same model in every role. For every turn we record the tool rounds and calls, whether the round limit or the budget forced the final answer, and the counts of malformed calls, repeated queries, evidence-budget hits and answers given without any search. These records yield a table of failures by type. % cite: Kapoor et al.; Kumar et al. 2026

\subsection{Cost and Latency Protocol}\label{sec:cost}

Cost is the number of input and output tokens of every model call multiplied by the list price of that model, read from the providers' price lists on 14 October 2026. The same rule applies to all models. It covers Qwen3.5-9B, which we self-host but price at its listed hosted price, the commercial models, which we reach through a subscription proxy, and the embedding and reranking models of retrieval. The reported cost is therefore what a deployment would pay, independent of how the experiments were run. Besides the cost per turn, we report the cost of 1{,}000 turns and the cost of pass, the cost per turn divided by the share of perfect answers. % cite: Erol et al. (cost of pass)

Latency comes from a separate run of each configuration at a concurrency of one, on the same machine and endpoints as the quality run, which may run four turns at a time. We report the median and 95th-percentile seconds per turn, and we split each turn into model, retrieval, reranking and decision seconds from the timestamps of its steps. Decision calls are timed at the client. Qwen3.5-9B runs on two NVIDIA T4 GPUs through a tunnel, while the commercial models answer through the proxy from Vietnam, so we compare latency only between architectures on the same model and never rank models by speed. Reasoning is disabled for Qwen3.5-9B and set to the lowest accepted level for Gemini 3.8 Flash. Claude Opus 5.5 does not allow reasoning to be disabled, so it runs at its minimal effort, and we report the reasoning tokens it actually spends.
